% ID: SEL-D0020-OA-E02
% ESTADO: cotejada. Segunda pasada visual completa y borrador compilado revisados.
% Fuente: Selectividad/matematicas_aplicadas/reserva b/Examen-Matemáticas Aplicadas a las Ciencias Sociales.pdf
% SHA-256: 37584a9604930b74008a6bda88ef3d61fa926d5e6a540ab9013569258f772256
% Página: 1; opción A; ejercicio original 2.

Los costes de producción de una empresa, en miles de euros, dependen de la cantidad de producto fabricada $x$, medida en toneladas, según la función $f(x)=30-9x+6x^2-x^3$. La capacidad máxima de producción es de 2 toneladas.
\begin{enumerate}[label=\alph*)]
\item \textbf{(1.25 puntos)} Obtenga los intervalos de crecimiento y decrecimiento de la función de costes de la empresa.
\item \textbf{(0.75 puntos)} Determine la cantidad que la empresa debe producir para minimizar los costes. \textquestiondown{}Cuál sería dicho coste mínimo?
\item \textbf{(0.5 puntos)} \textquestiondown{}Con qué producción la empresa tiene unos costes de producción máximos?
\end{enumerate}
